\subsection{Selección física como relación natural}
Fuera del centro electrónico, una especie no debe elegir por decreto una única celda. Las doce multisecciones no centrales son fibras estables bajo la inversión celular y no admiten, en general, una sección puntual equivariante. La realización correcta empieza por una relación

\[
\mathscr S\subset \mathcal P\times\mathcal R_{324},
\qquad
\mathscr S(X)=\{\gamma:(X,\gamma)\in\mathscr S\},
\]
cuyo valor puede ser un punto, una órbita o una fibra. Para el electrón, \(\mathscr S(e)=\{\gamma_e\}\) es de rango uno porque \((5,5)\) es el único punto fijo. Para color, sabor y generación, \(\mathscr S(X)\) conserva la órbita correspondiente hasta que un acoplamiento físico la comprime.

La relación es admisible si satisface simultáneamente:

\begin{enumerate}
\item depende sólo de representación, orientación, memoria, frontera, carga,
sabor, color y generación construidos antes del contraste;

\item conmuta con conjugación de partícula y antipartícula;
\item es equivariante respecto de la acción cromática y de la inversión celular;
\item es compatible con refinamiento, retorno nonádico y lectura dodecafásica;
\item conserva toda multiplicidad que no haya sido eliminada por un acoplamiento;
\item permanece invariante al permutar o reemplazar los valores externos de masa.
\end{enumerate}
El sexto punto es el control decisivo. Si alterar la columna experimental cambia \(\mathscr S(X)\), la construcción es retrospectiva. Si la relación permanece y cambia sólo el residual de contraste, la dirección causal está preservada.

\Needspace{7\baselineskip}
